\documentclass[14pt,a4paper,oneside]{extbook}
\usepackage[utf8]{vietnam}
\usepackage{amsmath,amssymb,mathrsfs,fancyhdr,enumerate,multirow,makecell}
\usepackage{tkz-euclide,tikz-3dplot,tikz,tkz-tab}
\usetikzlibrary{arrows.meta,calc,intersections,angles,quotes,backgrounds,shapes.geometric,patterns,through,shadows,positioning}
\usepackage{yhmath}
\usepackage[top=0.7cm, bottom=0.7cm, left=1.0cm, right=1.0cm, headheight=16pt, headsep=3mm, footskip=5mm]{geometry}
\usepackage[hidelinks,unicode]{hyperref}
\usepackage{qrcode}
\usepackage[many]{tcolorbox}
\usepackage{fontawesome}

\pagestyle{fancy}
\fancyhf{}
\fancyhead[L]{\footnotesize\bfseries\color{teal!85!black}\faBook~SỔ TAY KIẾN THỨC TOÁN 4}
\fancyhead[R]{\footnotesize\bfseries\color{blue!85!black}Lớp Toán Thầy Lĩnh -- loptoanthaylinh.vn}
\fancyfoot[C]{\footnotesize\bfseries\color{gray!80!black}-- Sổ tay kiến thức: Bài 72. Ôn tập thống kê và xác suất --}
\renewcommand{\headrulewidth}{0.4pt}

\begin{document}

% --- BANNER TIÊU ĐỀ BÀI HỌC ---
\noindent
\begin{tcolorbox}[
  enhanced,
  colback=teal!8!white,
  colframe=teal!80!black,
  arc=3mm, boxrule=1pt,
  left=3mm, right=3mm, top=2mm, bottom=2mm,
  halign=center
]
  {\Large\bfseries\color{teal!90!black}\faBookmark\ BÀI 72: ÔN TẬP THỐNG KÊ VÀ XÁC SUẤT}\\[1.5pt]
  {\footnotesize\itshape\color{black!75} Kĩ thuật đọc biểu đồ cột, tính tần số xuất hiện và dự đoán khả năng sự kiện trên 1 trang A4}
\end{tcolorbox}

\vspace{-2.5mm}

\noindent
\begin{minipage}[t]{0.49\textwidth}\vspace{0pt}
\begin{tcolorbox}[
  enhanced,
  colback=blue!4!white,
  colframe=blue!75!black,
  arc=2mm, boxrule=0.8pt,
  left=2mm, right=2mm, top=1.8mm, bottom=1.8mm,
  title={\footnotesize\bfseries\color{white}\faBarChart\ 1. Biểu đồ cột \& Thống kê}
]
{\footnotesize
\begin{itemize}\setlength{\itemsep}{1pt}\setlength{\leftskip}{-2.5mm}
  \item \textbf{Biểu đồ cột:} Trục ngang ghi đối tượng, trục đứng ghi số lượng kèm đơn vị.
  \item \textbf{Phân tích số liệu:}
  \begin{itemize}\setlength{\itemsep}{0.5pt}\setlength{\leftskip}{-2.5mm}
    \item Tìm giá trị lớn nhất (cột cao nhất).
    \item Tìm giá trị bé nhất (cột thấp nhất).
    \item Tính tổng số lượng của toàn biểu đồ.
    \item Tính số trung bình cộng của các cột.
  \end{itemize}
\end{itemize}
}
\end{tcolorbox}
\end{minipage}
\hfill
\begin{minipage}[t]{0.49\textwidth}\vspace{0pt}
\begin{tcolorbox}[
  enhanced,
  colback=teal!4!white,
  colframe=teal!75!black,
  arc=2mm, boxrule=0.8pt,
  left=2mm, right=2mm, top=1.8mm, bottom=1.8mm,
  title={\footnotesize\bfseries\color{white}\faQuestionCircle\ 2. Xác suất thực nghiệm}
]
{\footnotesize
\begin{itemize}\setlength{\itemsep}{1.5pt}\setlength{\leftskip}{-2.5mm}
  \item \textbf{Kiểm đếm:} Ghi chép số lần xảy ra sự kiện qua từng phép thử.
  \item \textbf{Quy luật tổng quát:}
  $$\text{Tổng số lần các sự kiện} = \text{Tổng số lần thử nghiệm}$$
  \item Đánh giá sự kiện xuất hiện nhiều nhất / ít nhất trong chuỗi thực nghiệm.
\end{itemize}
}
\end{tcolorbox}
\end{minipage}

\vspace{-2.5mm}
% --- KHỐI 3: SƠ ĐỒ TƯ DUY TIKZ ---
\noindent
\begin{tcolorbox}[
  enhanced,
  colback=orange!4!white,
  colframe=orange!80!black,
  arc=2mm, boxrule=0.8pt,
  left=2mm, right=2mm, top=1.5mm, bottom=1.5mm,
  title={\footnotesize\bfseries\color{white}\faSitemap\ SƠ ĐỒ TƯ DUY: TIẾN TRÌNH XỬ LÍ DỮ LIỆU THỐNG KÊ \& XÁC SUẤT}
]
\centering
\begin{tikzpicture}[
  box/.style={draw, rounded corners=2pt, align=center, font=\scriptsize\bfseries, inner sep=3pt, minimum height=1.05cm},
  arr/.style={-{Stealth[length=2.5mm]}, thick, shorten >=1.5pt, shorten <=1.5pt}
]
  \node[box, fill=purple!12, draw=purple!80!black, text width=3.2cm] (b1) at (0,0) {1. THU THẬP\\Khảo sát số liệu\\Thực nghiệm kiểm đếm};

  \node[box, fill=blue!10, draw=blue!80!black, text width=3.4cm] (b2) at (4.5,0) {2. TRỰC QUAN HÓA\\Lập bảng thống kê\\Vẽ biểu đồ cột};

  \node[box, fill=teal!10, draw=teal!80!black, text width=3.4cm] (b3) at (9.2,0) {3. PHÂN TÍCH\\Tìm Min, Max, TBC\\Đánh giá khả năng};

  \draw[arr, purple!80!black] (b1) -- (b2);
  \draw[arr, blue!80!black] (b2) -- (b3);
\end{tikzpicture}
\end{tcolorbox}

\vspace{-2.5mm}
% --- KHỐI 4: VÍ DỤ & MẸO TRÁNH BẪY ---
\noindent
\begin{minipage}[t]{0.49\textwidth}\vspace{0pt}
\begin{tcolorbox}[
  enhanced,
  colback=green!4!white,
  colframe=green!60!black,
  arc=2mm, boxrule=0.8pt,
  left=2mm, right=2mm, top=1.8mm, bottom=1.8mm,
  title={\footnotesize\bfseries\color{white}\faCheckCircle\ Ví dụ thực hành chuẩn mực}
]
{\footnotesize
\textbf{Ví dụ:} Bạn Mai quay tấm bìa 30 lần:
Mũi tên chỉ vào ô Xanh: 18 lần; ô Đỏ: 12 lần.
\begin{itemize}\setlength{\itemsep}{1pt}\setlength{\leftskip}{-2.5mm}
  \item Số lần xuất hiện ô Xanh \textbf{nhiều hơn} ô Đỏ: $18 - 12 = \mathbf{6}$ lần.
  \item Tổng số lần xuất hiện: $18 + 12 = \mathbf{30}$ lần.
\end{itemize}
}
\end{tcolorbox}
\end{minipage}
\hfill
\begin{minipage}[t]{0.49\textwidth}\vspace{0pt}
\begin{tcolorbox}[
  enhanced,
  colback=red!4!white,
  colframe=red!70!black,
  arc=2mm, boxrule=0.8pt,
  left=2mm, right=2mm, top=1.8mm, bottom=1.8mm,
  title={\footnotesize\bfseries\color{white}\faExclamationTriangle\ Mẹo tránh bẫy trắc nghiệm}
]
{\footnotesize
\begin{itemize}\setlength{\itemsep}{1.5pt}\setlength{\leftskip}{-2.5mm}
  \item \textbf{Bẫy vạch chia trên trục số:}\\
  Luôn quan sát vạch số đầu tiên xem mỗi khoảng cách tương ứng với $2, 5$ hay $10$ đơn vị.
  \item \textbf{Bẫy chia nhầm số phần tử:}\\
  Khi tính TBC của biểu đồ, phải chia cho số lượng các cột, không chia cho tổng số liệu.
\end{itemize}
}
\end{tcolorbox}
\end{minipage}

\end{document}
